\section*{Capítulo \ref{ch:g4:what-a-fire-needs} --- La combustión: lo que necesita un fuego}

\begin{solution}{exo:g4:what-a-fire-needs:1}
El \omterm{def:g4:what-a-fire-needs:fuel}{combustible}, el oxígeno (del \omterm{def:g4:air-a-mixture-of-gases:air}{aire}) y el calor (para que empiece el
fuego).
\end{solution}

\begin{solution}{exo:g4:what-a-fire-needs:2}
La madera, la cera de vela, la gasolina y el papel son \omterm{def:g4:what-a-fire-needs:fuel}{combustibles}. La
arena, el vidrio y el agua no \omterm{def:g4:what-a-fire-needs:burning}{arden}.
\end{solution}

\begin{solution}{exo:g4:what-a-fire-needs:3}
El oxígeno: la manta impide que el \omterm{def:g4:air-a-mixture-of-gases:air}{aire} llegue al fuego.
\end{solution}

\begin{solution}{exo:g4:what-a-fire-needs:4}
El calor: el agua enfría tanto la madera que \omterm{def:g4:what-a-fire-needs:burning}{arde} que ya no puede
seguir \omterm{def:g4:what-a-fire-needs:burning}{ardiendo}.
\end{solution}

\begin{solution}{exo:g4:what-a-fire-needs:5}
Tres cualesquiera de estas: ceniza, humo, dióxido de carbono, vapor de
agua.
\end{solution}

\begin{solution}{exo:g4:what-a-fire-needs:6}
Al soplar llega más \omterm{def:g4:air-a-mixture-of-gases:air}{aire}, y por tanto más oxígeno, al \omterm{def:g4:what-a-fire-needs:fuel}{combustible}
caliente, que vuelve a \omterm{def:g4:what-a-fire-needs:burning}{arder}.
\end{solution}

\begin{solution}{exo:g4:what-a-fire-needs:7}
Agua: la cera que \omterm{def:g4:what-a-fire-needs:burning}{arde} produce vapor de agua, que forma gotitas sobre el
vidrio frío.
\end{solution}

\begin{solution}{exo:g4:what-a-fire-needs:8}
Salir, agachándote por debajo del humo y cerrando las puertas al pasar,
y después llamar a los bomberos desde fuera.
\end{solution}

\begin{solution}{exo:g4:what-a-fire-needs:9}
El lado del \omterm{def:g4:what-a-fire-needs:fuel}{combustible}: con el gas cerrado, al fuego ya no le queda
nada que quemar.
\end{solution}

\begin{solution}{exo:g4:what-a-fire-needs:10}
El \omterm{def:g4:what-a-fire-needs:fuel}{combustible}: como en la franja no hay árboles, al fuego se le acaba
el \omterm{def:g4:what-a-fire-needs:fuel}{combustible} cuando llega a ella.
\end{solution}

\begin{solution}{exo:g4:what-a-fire-needs:11}
No. La mayor parte de la madera, junto con el oxígeno del \omterm{def:g4:air-a-mixture-of-gases:air}{aire}, se ha
convertido en sustancias nuevas que son gases (dióxido de carbono,
vapor de agua y humo) y que se han ido al \omterm{def:g4:air-a-mixture-of-gases:air}{aire}. Solo queda la ceniza.
\end{solution}
